\section{Coupled force and formation planning}
\subsection{Aperture-conditioned motion references}
The inputs are the reconstructed surface, calibrated attachments, known obstacle boxes, fitted inertial parameters and a prescribed smooth payload reference family. The continuous optimization variables are the three components of every cable force at each time node. Duration and lowering are selected by ordered finite candidate search. Payload position and attitude follow the selected reference; aircraft endpoints are derived from optimized forces and cable stretch. Thus force allocation and support formation adapt to the measured load, while the obstacle-crossing attitude sequence is supplied by the reference family.

The load reference is represented by smooth translation and attitude functions over a selected duration $D$. The basic experiment translates the payload by 8 m while moving toward a terminal roll, pitch and yaw of $25^\circ$, $10^\circ$ and $30^\circ$. The reference includes a smooth lowering window under the beam. Candidate durations are 20, 22, 25 and 30 s; candidate lowering values are evaluated in a fixed order.

For sequential apertures, a piecewise quintic centerline connects their approach and exit neighborhoods. Smooth activation functions blend the associated attitude references. In the large-angle course the sequence is horizontal, $90^\circ$ roll at the narrow aperture, horizontal with $90^\circ$ yaw under the ceiling, $-85^\circ$ roll at the tilted aperture, and return to the nominal attitude. The references vary during forward motion. Their activation regions extend across the occupied aperture neighborhood, so a load transition is not released while supporting aircraft still occupy the frame. Duration and lowering candidates are evaluated using the same force and endpoint checks.

\subsection{Smooth conic cable-force optimization}
At time nodes $t_j$, let $f_j\in\R^{3N}$ contain the load cable forces, and let $f_j^0$ be a preferred outward preload. For a chosen motion candidate, the force planner minimizes
\begin{equation}
 \begin{split}
 \mathcal{J}(f)=&\frac{0.02}{2}\norm{f-f^0}^{2}
 +\frac{0.04}{2}\norm{D_1f}^{2}\\
 &+\frac{0.20}{2}\norm{D_2f}^{2}
 +\frac{0.04}{2}\norm{D_3f}^{2},
 \end{split}
 \label{eq:objective}
\end{equation}
where $D_1,D_2,D_3$ are time-scaled first, second and third finite-difference operators. Force equilibrium is imposed at every node:
\begin{equation}
 G(R_r(t_j))f_j=Y_d(t_j)\hat\theta.
 \label{eq:equilibrium}
\end{equation}
The formulation includes upward force, tension and nominal thrust cones,
\begin{equation}
 f_{ij,z}\geq0.55,\quad \norm{f_{ij}}\leq8.2,\quad
 \norm{f_{ij}+m_ug_0 e_z}\leq29.5,
 \label{eq:cones}
\end{equation}
with forces in newtons, $e_z=[0,0,1]^T$, and $g_0=9.81$ m/s$^2$ denoting the positive gravity magnitude. In the thrust cone, $m_ug_0e_z=-m_u g$ with the vector convention of Section 2. The 29.5 N cone limits nominal total thrust; four 7.5 N motor limits permit 30 N total before moment allocation. Dynamic endpoint verification applies a tighter 28 N nominal limit. Componentwise limits of 8 N/s and 25 N/s$^2$ are placed on the first and second differences. Near-zero first differences at the endpoints reduce force transients when joining the preparation and transport phases. The initial force is projected onto the first wrench equality to maintain continuity.

An outward preload prevents the formation from collapsing above the object. For large roll angles, horizontal aircraft sectors are defined using a horizontal projection of the payload longitudinal axis and its orthogonal horizontal direction. The sectors retain the original attachment ordering when the attachment plane becomes nearly vertical. Side-force conditions are relaxed with increasing roll to allow the force couples required by the upper and lower attachments. The subsequent endpoint and cable checks determine acceptance.

\subsection{Aperture constraints and compliant endpoints}
For an aperture halfspace $n^Tx_i\leq\beta$ and known attachment $a_i$, define the endpoint allowance $c_i=\beta-n^Ta_i>0$. A conservative linear force condition used in the conic problem is
\begin{equation}
 \left(n-\frac{c_i}{\ell_{\max}}v_i\right)^Tf_i\leq0,
 \label{eq:halfspace}
\end{equation}
where $v_i$ is a prescribed unit direction and $\ell_{\max}=2.05$ m bounds the planned cable length. Since $v_i^Tf_i\leq\norm{f_i}$, this condition implies $n^Tf_i\leq(c_i/\ell_{\max})\norm{f_i}$. Together with $\ell_i\leq\ell_{\max}$, it bounds the nonlinear endpoint displacement toward the halfspace. Height limits are expressed by related outward/upward force inequalities. Full nonlinear geometry is checked after solving.

For a smooth force magnitude $T_i(t)=\norm{f_i(t)}$, stretch is computed from
\begin{equation}
 d_c\dot\delta_i+k_c\delta_i=T_i(t),\qquad
 x_{i,r}=p_r+R_rb_i+(\ell_0+\delta_i)\frac{f_i}{T_i}.
 \label{eq:stretch}
\end{equation}
The stretch dynamics are propagated at 5 ms resolution, and clamped cubic splines supply the endpoint position, velocity and acceleration. This step retains axial damping instead of instantaneously changing cable length whenever tension changes. Aircraft acceleration, dynamic thrust and tilt are evaluated on the resulting endpoint curves, supplementing the nominal cones in Eq.~\eqref{eq:cones}.

\subsection{Separation during large-angle reorientation}
Near-vertical payload orientations can bring different endpoint directions close together even when the attachments themselves are well separated. We therefore add sequential linearized endpoint separation. At a preceding force solution $\bar f_i$, the static elastic map in Eq.~\eqref{eq:endpointmap} has Jacobian
\begin{equation}
 J_i=\frac{\ell_0}{\norm{\bar f_i}}\left(I-\bar q_i\bar q_i^T\right)+\frac{1}{k_c}I.
 \label{eq:jacobian}
\end{equation}
For the previous endpoint difference direction $n_{i\ell}$, impose
\begin{equation}
 n_{i\ell}^{T}\left[\bar x_i+J_i(f_i-\bar f_i)
 -\bar x_\ell-J_\ell(f_\ell-\bar f_\ell)\right]\geq d_{\rm lin}.
 \label{eq:separation}
\end{equation}
Three successive conic updates use $d_{\rm lin}=0.72$ m. The final damped endpoint trajectory is then evaluated nonlinearly. Its acceptance distance is 0.52 m; the linearization distance is an optimization reference rather than a guaranteed executed separation. This distinction matters because damping, spline interpolation and tracking alter the static linearization geometry.

Feasibility is evaluated at three distinct levels. The conic solve enforces nodal wrench equalities and force cones. Candidate verification evaluates damped, interpolated endpoints and complete geometric occupancy. Closed-loop execution is assessed using the measured tracking, loading and actuation limits in Section 5. The 0.72 m separation linearization reference encourages spacing during optimization; the nonlinear candidate gate is 0.52 m. Execution reports aircraft center distance separately from clearance after subtracting the collision envelope. Preparation transients are also included in execution tension extrema; they are not subject to the transport-only force reference bound.

The course uses 121 force nodes. Final evaluation samples the full acquired surface, aircraft envelopes and complete cable segments at 20 ms intervals, with actuation evaluation on the finer endpoint curve. A candidate is accepted only after its nonlinear force, geometry, speed and acceleration checks pass. Accepted course references satisfy aircraft speed below 2.5 m/s, reference acceleration below 2.75 m/s$^2$, tension between 1 and 8.25 N, dynamic thrust below 28 N and tilt below $37^\circ$. The single-beam reference acceleration budget is 2.25 m/s$^2$. These budgets leave room for tracking corrections in the execution loop.

\subsection{Closed-loop execution}
At each control instant, a delayed registered load pose is propagated using its available velocity estimate and combined with current aircraft telemetry. A low-pass payload pose correction shifts the endpoint references. Wider corner mounts use a lower visual pose gain to keep delayed rotational corrections below the coupled load-swing response. The aircraft command is
\begin{equation}
 u_i=m_u\left(\ddot x_{i,r}+K_p(x_{i,r}+\Delta_i-x_i)
      +K_v(\dot x_{i,r}-\dot x_i)-g\right)+f_i+d_v\dot x_i,
 \label{eq:controller}
\end{equation}
where $K_p=18I$ s$^{-2}$ and $K_v=7I$ s$^{-1}$, $\Delta_i$ is the filtered pose correction, and requested acceleration is capped at 4.5 m/s$^2$. A geometric attitude controller~\cite{lee2010} aligns each aircraft's thrust direction with this command, followed by allocation to four filtered motor inputs.

Image registration runs in its own worker. Identification and whole-segment planning run in a separate process; the aircraft remain under preparation/hover control while these computations finish. After the candidate is ready, the transport clock starts and the 50 Hz feedback loop follows its reference. This arrangement assigns explicit computation and observation times to the complete sensing--planning--execution chain.
